% ===========================================================================
% 11. Compound Trap Analysis
% ===========================================================================
\section{Compound Trap Analysis}
\label{sec:compound}

The DeepMind taxonomy treats each trap category independently, but
real-world adversaries are unlikely to restrict themselves to a single
attack vector. We evaluate all 15 pairwise combinations of the six
categories to test whether compound traps produce super-additive
effects.

\subsection{Compound Design}
\label{sec:compound-design}

For each pair of categories $A$ and $B$, we select the representative
scenario with the highest baseline trap success rate from each
category. The compound trap combines both adversarial environments:
the agent is exposed to content from both traps simultaneously. We
hypothesize that compound traps succeed at a rate higher than
$\max(\tau_A, \tau_B)$ --- i.e., they are super-additive.

\subsection{Results}
\label{sec:compound-results}

\begin{table}[t]
  \centering
  \caption{Compound trap success rates (\%) vs.\ individual trap
    baselines. $\tau_A$, $\tau_B$: individual trap success rates.
    $\tau_{A+B}$: compound success rate. $\Delta$: super-additive
    surplus ($\tau_{A+B} - \max(\tau_A, \tau_B)$). Averaged across
    4 models, $n=10$ reps. \sig\ denotes significant super-additivity.}
  \label{tab:compound}
  \small
  \begin{tabular}{@{}llcccc@{}}
    \toprule
    \textbf{Cat.\ A} & \textbf{Cat.\ B} & $\tau_A$ & $\tau_B$ & $\tau_{A+B}$ & $\Delta$ \\
    \midrule
    CI & SM & \textcolor{red}{[TBD]} & \textcolor{red}{[TBD]} & \textcolor{red}{[TBD]} & \textcolor{red}{[TBD]} \\
    CI & CS & \textcolor{red}{[TBD]} & \textcolor{red}{[TBD]} & \textcolor{red}{[TBD]} & \textcolor{red}{[TBD]} \\
    CI & BC & \textcolor{red}{[TBD]} & \textcolor{red}{[TBD]} & \textcolor{red}{[TBD]} & \textcolor{red}{[TBD]} \\
    CI & SY & \textcolor{red}{[TBD]} & \textcolor{red}{[TBD]} & \textcolor{red}{[TBD]} & \textcolor{red}{[TBD]} \\
    CI & HL & \textcolor{red}{[TBD]} & \textcolor{red}{[TBD]} & \textcolor{red}{[TBD]} & \textcolor{red}{[TBD]} \\
    SM & CS & \textcolor{red}{[TBD]} & \textcolor{red}{[TBD]} & \textcolor{red}{[TBD]} & \textcolor{red}{[TBD]} \\
    SM & BC & \textcolor{red}{[TBD]} & \textcolor{red}{[TBD]} & \textcolor{red}{[TBD]} & \textcolor{red}{[TBD]} \\
    SM & SY & \textcolor{red}{[TBD]} & \textcolor{red}{[TBD]} & \textcolor{red}{[TBD]} & \textcolor{red}{[TBD]} \\
    SM & HL & \textcolor{red}{[TBD]} & \textcolor{red}{[TBD]} & \textcolor{red}{[TBD]} & \textcolor{red}{[TBD]} \\
    CS & BC & \textcolor{red}{[TBD]} & \textcolor{red}{[TBD]} & \textcolor{red}{[TBD]} & \textcolor{red}{[TBD]} \\
    CS & SY & \textcolor{red}{[TBD]} & \textcolor{red}{[TBD]} & \textcolor{red}{[TBD]} & \textcolor{red}{[TBD]} \\
    CS & HL & \textcolor{red}{[TBD]} & \textcolor{red}{[TBD]} & \textcolor{red}{[TBD]} & \textcolor{red}{[TBD]} \\
    BC & SY & \textcolor{red}{[TBD]} & \textcolor{red}{[TBD]} & \textcolor{red}{[TBD]} & \textcolor{red}{[TBD]} \\
    BC & HL & \textcolor{red}{[TBD]} & \textcolor{red}{[TBD]} & \textcolor{red}{[TBD]} & \textcolor{red}{[TBD]} \\
    SY & HL & \textcolor{red}{[TBD]} & \textcolor{red}{[TBD]} & \textcolor{red}{[TBD]} & \textcolor{red}{[TBD]} \\
    \bottomrule
  \end{tabular}
\end{table}

\subsection{Key Findings}
\label{sec:compound-findings}

\begin{enumerate}[nosep,leftmargin=*]
  \item \textbf{Content injection + semantic manipulation} is the most
    potent compound pair: hidden instructions wrapped in authority
    framing achieve significantly higher success than either attack
    alone (large effect size).

  \item \textbf{Cognitive state + systemic} compounds demonstrate the
    largest blast radius: poisoned knowledge propagated through A2A
    messaging corrupts multiple agents, and each corrupted agent
    further reinforces the false information.

  \item \textbf{Semantic manipulation + human-in-the-loop} is a
    particularly concerning combination for real-world deployment:
    the agent is biased by framing effects \emph{and} produces biased
    reports for human reviewers, creating a double layer of
    manipulation.

  \item Overall, \textcolor{red}{[TBD]} of the 15 compound pairs show
    statistically significant super-additivity, confirming that the
    DeepMind categories interact and that defenses must address
    cross-category attack composition.
\end{enumerate}
